\documentclass[10pt]{article}

% COMP 2300 Divisibility Core companion: student-facing reference.
\usepackage[letterpaper,margin=0.56in,includehead,includefoot,
  headheight=15pt,headsep=7pt,footskip=20pt]{geometry}
\usepackage{amsmath,amssymb,mathtools}
\usepackage{microtype}
\usepackage[dvipsnames,svgnames]{xcolor}
\usepackage{multicol}
\usepackage{enumitem}
\usepackage{tabularx,booktabs,array}
\usepackage[most]{tcolorbox}
\usepackage{fancyhdr}
\usepackage{titlesec}
\usepackage{needspace}
\usepackage[hidelinks]{hyperref}

\definecolor{DivNavy}{HTML}{18364A}
\definecolor{DivTeal}{HTML}{176B6B}
\definecolor{DivGold}{HTML}{9A6500}
\definecolor{DivRed}{HTML}{963D38}
\definecolor{DivBlueTint}{HTML}{EDF4F8}
\definecolor{DivTealTint}{HTML}{EDF8F6}
\definecolor{DivGoldTint}{HTML}{FFF7E5}
\definecolor{DivRedTint}{HTML}{FCEFED}

\setlength{\parindent}{0pt}
\setlength{\parskip}{3pt}
\setlength{\columnsep}{0.25in}
\setlength{\columnseprule}{0.25pt}
\setlength{\emergencystretch}{1em}
\setlength{\abovedisplayskip}{4pt}
\setlength{\belowdisplayskip}{4pt}
\setlength{\abovedisplayshortskip}{2pt}
\setlength{\belowdisplayshortskip}{2pt}
\setlist{leftmargin=*,itemsep=1.5pt,topsep=2pt,parsep=0pt,partopsep=0pt}
\raggedbottom

\titleformat{\section}{\Large\bfseries\color{DivNavy}}{}{0pt}{}
  [\vspace{-3pt}\color{DivNavy}\rule{\linewidth}{0.7pt}]
\titleformat{\subsection}{\large\bfseries\color{DivTeal}}{}{0pt}{}
\titlespacing*{\section}{0pt}{4pt}{5pt}
\titlespacing*{\subsection}{0pt}{5pt}{2pt}

\pagestyle{fancy}
\fancyhf{}
\fancyhead[L]{\textbf{COMP 2300}}
\fancyhead[C]{\textbf{Divisibility Core Companion}}
\fancyhead[R]{Definitions and Worked Examples}
\fancyfoot[L]{Core sections A--E}
\fancyfoot[C]{\thepage}
\fancyfoot[R]{Name the definition you use}
\renewcommand{\headrulewidth}{0.5pt}
\renewcommand{\footrulewidth}{0.3pt}
\fancypagestyle{plain}{%
  \fancyhf{}%
  \fancyhead[L]{\textbf{COMP 2300}}%
  \fancyhead[C]{\textbf{Divisibility Core Companion}}%
  \fancyhead[R]{Definitions and Worked Examples}%
  \fancyfoot[L]{Core sections A--E}%
  \fancyfoot[C]{\thepage}%
  \fancyfoot[R]{Name the definition you use}%
  \renewcommand{\headrulewidth}{0.5pt}%
  \renewcommand{\footrulewidth}{0.3pt}%
}
\makeatletter
\let\ps@plain\ps@fancy
\let\ps@empty\ps@fancy
\makeatother

\tcbset{enhanced,before skip=4pt,after skip=4pt,boxrule=0.6pt,arc=1.4mm,
  left=5pt,right=5pt,top=4pt,bottom=4pt,fonttitle=\bfseries}
\newtcolorbox{definitionbox}[1]{title={#1},colback=DivBlueTint,
  colframe=DivNavy,colbacktitle=DivNavy,coltitle=white}
\newtcolorbox{methodbox}[1]{title={#1},colback=DivTealTint,
  colframe=DivTeal,colbacktitle=DivTeal,coltitle=white}
\newtcolorbox{workedbox}[1]{title={#1},colback=white,
  colframe=DivNavy,colbacktitle=DivBlueTint,coltitle=DivNavy}
\newtcolorbox{warningbox}[1]{title={#1},colback=DivRedTint,
  colframe=DivRed,colbacktitle=DivRed,coltitle=white}
\newtcolorbox{checklistbox}[1]{title={#1},colback=DivGoldTint,
  colframe=DivGold,colbacktitle=DivGold,coltitle=white}
\newcommand{\Z}{\mathbb Z}
\newcommand{\term}[1]{\textbf{\color{DivNavy}#1}}
\newcommand{\equivmod}[1]{\mathrel{\equiv_{#1}}}

\begin{document}
\makeatletter
\let\ps@plain\ps@fancy
\let\ps@empty\ps@fancy
\makeatother
\pagestyle{fancy}

\section*{A. Division, Remainders, and Congruence}

\begin{multicols}{2}
\begin{definitionbox}{Divisibility}
For integers $d$ and $n$,
\[
 d\mid n\quad\Longleftrightarrow\quad
 \text{there is a }k\in\Z\text{ such that }n=dk.
\]
Here $\Z$ is the set of integers, $d$ is a \term{divisor} (or factor), $n$
is a \term{multiple} of $d$, and $k$ is an integer \term{witness}. To prove
$d\nmid n$, show that no integer witness can exist.
\end{definitionbox}

\begin{definitionbox}{Division Algorithm}
For every $a\in\Z$ and positive integer $b$, there are unique integers $q,r$
such that
\[
 a=bq+r,\qquad 0\le r<b.
\]
The integer $a$ is the \term{dividend}, $b$ the \term{divisor}, $q$ the
\term{quotient}, and $r$ the \term{remainder}. We write
$q=a\,\mathrm{div}\,b$ and $r=a\bmod b$. The bound on $r$ makes the answer
canonical, even when $a$ is negative.
\end{definitionbox}

\begin{definitionbox}{Congruence}
For $m\ge2$,
\[
 x\equiv y\pmod m
 \quad\Longleftrightarrow\quad m\mid(x-y).
\]
The positive integer $m$ is the \term{modulus}; we may also write
$x\equivmod{m}y$. Equivalently, $x=y+mk$ for some $k\in\Z$, or $x$ and $y$
have the same least nonnegative remainder modulo $m$.
\end{definitionbox}

\begin{methodbox}{One relationship, three views}
If $a=bq+r$ with $0\le r<b$, then
\[
 a\bmod b=r
 \quad\text{and}\quad
 a\equiv r\pmod b.
\]
The equation gives the quotient and remainder; congruence records only the
remainder class.
\end{methodbox}

\begin{methodbox}{Basic divisibility laws}
For integers of the indicated form:
\[
\begin{array}{ll}
d\mid a,\ d\mid b &\Longrightarrow d\mid(ma+nb),\quad m,n\in\Z,\\
d\mid a &\Longrightarrow d\mid ac,\quad c\in\Z,\\
d\mid a,\ a\mid b &\Longrightarrow d\mid b.
\end{array}
\]
Each follows by substituting integer witnesses and factoring out $d$.
\end{methodbox}

\columnbreak

\begin{workedbox}{Worked example: a positive dividend}
Find all three views for $263$ divided by $17$.
\textbf{Choose $q$:} $17(15)=255$ is the largest multiple not exceeding $263$.
\textbf{Compute $r$:} $263-255=8$, so
\[
263=17(15)+8,\qquad 0\le8<17.
\]
The bound $0\le8<17$ makes this canonical. Therefore
\[
\begin{aligned}
263\,\mathrm{div}\,17&=15, & 263\bmod17&=8,\\
263&\equiv8\pmod{17}.
\end{aligned}
\]
\textbf{Congruence check:}
$263-8=255=17(15)$, so $17\mid(263-8)$.
\end{workedbox}

\begin{workedbox}{Worked example: a negative dividend}
Find the canonical quotient and remainder for $-52$ divided by $9$.
\textbf{Reject the first try:} $-52=9(-5)-7$ has a negative remainder.
\textbf{Adjust:} decrease the quotient by $1$ and add $9$ to the remainder:
\[
-52=9(-6)+2,\qquad 0\le2<9.
\]
Now $0\le2<9$. Thus
\[
\begin{aligned}
-52\,\mathrm{div}\,9&=-6, & -52\bmod9&=2,\\
-52&\equiv2\pmod9.
\end{aligned}
\]
\end{workedbox}

\begin{warningbox}{Common mistake}
Do not truncate a negative quotient toward zero. Choose the unique $q$ that
makes $0\le r<b$. A negative number may be congruent to a negative residue,
but $a\bmod b$ is the least nonnegative one.
\end{warningbox}

\begin{checklistbox}{Can you do this from the definitions?}
Given $a$ and $b>0$, find $q,r$; verify the remainder bound; state
$a\,\mathrm{div}\,b$, $a\bmod b$, and a congruence; justify the congruence by
showing that $b$ divides a difference.
\end{checklistbox}

\begin{workedbox}{Recover a number; test a congruence}
\textbf{Recover:} insert $q=-4,r=9$ into $a=13q+r$ to get
$a=13(-4)+9=-43$; $0\le9<13$ verifies the remainder.
\textbf{Test:} subtract in $347\equiv-13\pmod{24}$:
$347-(-13)=360=24(15)$, so the congruence is true. Finally,
$347=24(14)+11$, so the least nonnegative residue is $11$.
\end{workedbox}
\end{multicols}

\thispagestyle{fancy}
\clearpage
\pagestyle{fancy}
\thispagestyle{fancy}
\section*{B. Base-$b$ Representations}
\thispagestyle{fancy}

\begingroup
\tcbset{before skip=3pt,after skip=3pt}
\begin{multicols}{2}\small
\begin{definitionbox}{Base-$b$ expansion}
For $b\ge2$, the numeral $(a_ka_{k-1}\cdots a_0)_b$ represents
\[
 a_kb^k+a_{k-1}b^{k-1}+\cdots+a_1b+a_0,
 \qquad 0\le a_i<b.
\]
Binary uses digits $0,1$; octal uses $0$ through $7$; hexadecimal uses
$0$ through $9$ and $A=10,\ldots,F=15$.
\end{definitionbox}

\begin{methodbox}{Place value, shifting, and digits}
\[
(10)_b=b,\qquad (100)_b=b^2,\qquad
(10)_b(10)_b=(100)_b.
\]
More generally,
\[
\bigl((10)_b\bigr)^k
=\bigl(1\underbrace{00\cdots0}_{k\text{ zeros}}\bigr)_b=b^k.
\]
If $N=(a_k\cdots a_0)_b$ and $0\le r<b$, then
\[
bN=(a_k\cdots a_00)_b,
\qquad bN+r=(a_k\cdots a_0r)_b.
\]
Equivalently,
\[
(a_k\cdots a_1a_0)_b=b(a_k\cdots a_1)_b+a_0.
\]
Thus multiplying by $b$ shifts left, while dividing by $b$ removes the final
digit: $\mathrm{div}\,b$ gives the remaining prefix and $\bmod b$ gives the
removed digit.
\end{methodbox}

\begin{methodbox}{Decimal to another base}
Repeatedly divide by the new base. Record each remainder. The digits are the
remainders read from \emph{bottom to top}. Verify by expanding in powers of the
base. Each line $n=bq+r$ reverses one shift-and-append step. In hexadecimal,
remainders $10$--$15$ become $A$--$F$; for $1000_{10}$ the successive
remainders are $8,E,3$, so $1000_{10}=(3E8)_{16}$.
\end{methodbox}

\begin{workedbox}{Worked example: $202_{10}$ to binary}
\textbf{Divide and record remainders:}
\[
\begin{array}{rcl@{\qquad}rcl}
202&=&2(101)+0 & 101&=&2(50)+1\\
50&=&2(25)+0 & 25&=&2(12)+1\\
12&=&2(6)+0 & 6&=&2(3)+0\\
3&=&2(1)+1 & 1&=&2(0)+1
\end{array}
\]
\textbf{Read upward:} the remainders in reverse chronological order give
$202_{10}=11001010_2$. Rebuilding by shifting and
appending verifies the digit order:
\[
\begin{aligned}
(1)_2&\to(11)_2\to(110)_2\\
&\to(1100)_2\to(11001)_2\to(110010)_2\\
&\to(1100101)_2\to(11001010)_2.
\end{aligned}
\]
\textbf{Place-value check:} $128+64+8+2=202$.
\end{workedbox}

\columnbreak

\begin{workedbox}{Worked example: evaluate a numeral}
Multiply each digit by the matching power of $5$:
\[
(2301)_5=2(5^3)+3(5^2)+0(5)+1=326.
\]
Thus the value is $326_{10}$. The subscript names the base; it is not a
multiplier.
\end{workedbox}

\begin{methodbox}{Binary as a bridge}
One hexadecimal digit equals exactly four binary bits; one octal digit equals
exactly three. Pad with leading zeros, then regroup from the \emph{right}. For
example,
\[
\begin{aligned}
(725)_8&=(111\ 010\ 101)_2=469_{10},\\
(101\ 101\ 011)_2&=(553)_8,\\
(0001\ 0110\ 1011)_2&=(16B)_{16}.
\end{aligned}
\]
\end{methodbox}

\begin{workedbox}{Worked example: hexadecimal to binary and octal}
Convert $(2D7)_{16}$.
\textbf{Hex to binary:} replace each digit with four bits:
\[
2=0010_2,\qquad D=1101_2,\qquad7=0111_2,
\]
so
\[
(2D7)_{16}=(0010\ 1101\ 0111)_2.
\]
\textbf{Binary to octal:} regroup from the right in threes,
\[
001\ 011\ 010\ 111
\quad\Longrightarrow\quad (1327)_8.
\]
\textbf{Decimal check:}
\[
2(16^2)+13(16)+7=727.
\]
Leading zeros change grouping, not value.
\end{workedbox}

\begin{workedbox}{Worked example: the last two base-$b$ digits}
Let $n=(25314)_7$. Then
\[
n=2(7^4)+5(7^3)+3(7^2)+1(7)+4.
\]
\textbf{Discard higher places:} every term containing $7^2$ is congruent to
$0$ modulo $7^2$. Therefore
\[
n\bmod7^2=1(7)+4=11.
\]
\textbf{Interpret:} the last two base-seven digits are $(14)_7$, whose decimal
value is $1(7)+4=11$, not $14$.
\end{workedbox}

\begin{checklistbox}{Conversion checks}
Digits must be legal for the base. Binary/hex groups have four bits; binary/octal
groups have three. Expand the final numeral to verify it represents the same
integer.
\end{checklistbox}
\end{multicols}
\endgroup

\thispagestyle{fancy}
\clearpage
\pagestyle{fancy}
\thispagestyle{fancy}
\section*{C--D. Modular Arithmetic, GCD, and B\'ezout}
\thispagestyle{fancy}

\begingroup
\tcbset{before skip=0.5pt,after skip=0.5pt}
\setlength{\abovedisplayskip}{3pt}
\setlength{\belowdisplayskip}{3pt}
\begin{multicols}{2}\small
\begin{definitionbox}{Arithmetic with congruences}
If $a\equiv b\pmod m$ and $c\equiv d\pmod m$, then
\[
a+c\equiv b+d\pmod m,
\qquad ac\equiv bd\pmod m.
\]
In particular, $a^n\equiv b^n\pmod m$ for every positive integer $n$.
The exact remainder identities are
\[
\begin{aligned}
(a+c)\bmod m&=((a\bmod m)+(c\bmod m))\bmod m,\\
ac\bmod m&=((a\bmod m)(c\bmod m))\bmod m.
\end{aligned}
\]
Reduce early; report a least nonnegative residue for $x\bmod m$.
\end{definitionbox}

\begin{workedbox}{Worked example: reduce early}
Compute the least nonnegative residue of $(-47)\cdot38+19$ modulo $12$.
\textbf{Reduce each input:}
\[
-47\equiv1,\quad38\equiv2,\quad19\equiv7\pmod{12},
\]
\textbf{Operate on the small representatives:}
\[
(-47)\cdot38+19\equiv1(2)+7\equiv9\pmod{12}.
\]
Since $0\le9<12$, $((-47)\cdot38+19)\bmod12=9$.
\end{workedbox}

\begin{methodbox}{Powers modulo $m$: cycles and squaring}
\textbf{Cycle method.} Compute residues until they repeat. Since
$3,9,7,1$ repeat for powers of $3$ modulo $10$, and $47\equiv3\pmod4$,
\[
3^{47}\equiv3^3\equiv7\pmod{10}.
\]
\textbf{Repeated squaring.} Square and reduce at every step. Modulo $13$,
\[
5^2\equiv-1,\qquad5^4\equiv1,
\]
so all powers $5^{4j}\equiv1$. Because
$117=64+32+16+4+1$,
\[
5^{117}\equiv1\cdot1\cdot1\cdot1\cdot5\equiv5\pmod{13}.
\]
State the cycle length or the exponent decomposition you used.
\end{methodbox}

\columnbreak

\begin{definitionbox}{GCD and relative primality}
For integers $a,b$, not both zero, $\gcd(a,b)$ is the greatest positive
integer dividing both. The numbers are \term{relatively prime} when
$\gcd(a,b)=1$. The numbers themselves need not be prime.
\end{definitionbox}

\begin{methodbox}{Euclidean algorithm}
Put the larger positive input first, then repeatedly use
\[
a=bq+r,\qquad0\le r<b,
\]
then replace $(a,b)$ by $(b,r)$. Stop when the remainder is $0$. The last
\emph{nonzero} remainder is the gcd. Indeed, from $a=bq+r$, the common
divisors of $a,b$ are exactly those of $b,r=a-bq$, so
$\gcd(a,b)=\gcd(b,r)$. Swap the inputs first if necessary.
\end{methodbox}

\begin{workedbox}{Euclidean algorithm and B\'ezout coefficients}
Find $\gcd(246,102)$ and express it as $246s+102t$.
\textbf{Step 1: run the Euclidean algorithm.}
\[
\begin{aligned}
246&=2(102)+42,\\
102&=2(42)+18,\\
42&=2(18)+6,\\
18&=3(6)+0.
\end{aligned}
\]
The last nonzero remainder is $6$, so it is the gcd. \textbf{Step 2:
back-substitute.} B\'ezout's identity says that
$\gcd(a,b)=sa+tb$ for some $s,t\in\Z$. Back-substitute:
\[
\begin{aligned}
6&=42-2(18)\\
 &=42-2\bigl(102-2(42)\bigr)\\
 &=5(42)-2(102)\\
 &=5\bigl(246-2(102)\bigr)-2(102)\\
 &=5(246)-12(102).
\end{aligned}
\]
Thus $s=5$ and $t=-12$. \textbf{Step 3: verify both conclusions:}
\[
6\mid246,\quad6\mid102,\qquad5(246)-12(102)=6.
\]
\end{workedbox}

\begin{methodbox}{B\'ezout and relative-prime checks}
Back-substitute through every displayed Euclidean equation; preserve the
original input order in the final identity. Coefficients are not unique, so
verify them by substitution. For example,
$16=2^4$ and $45=3^2\cdot5$ share no prime factor, so
$\gcd(16,45)=1$: composite numbers can be relatively prime.
\end{methodbox}
\end{multicols}
\endgroup

\thispagestyle{fancy}
\clearpage
\pagestyle{fancy}
\thispagestyle{fancy}
\section*{E. Definition-Driven Proofs and Quiz Bridges}
\thispagestyle{fancy}

\begingroup
\tcbset{before skip=3pt,after skip=3pt}
\begin{multicols}{2}\small
\begin{definitionbox}{Even, odd, prime, and composite}
An integer is \term{even} if it is $2k$, and \term{odd} if it is $2k+1$,
for some $k\in\Z$. A positive integer $p>1$ is \term{prime} if its only
positive divisors are $1$ and $p$; an integer greater than $1$ that is not
prime is \term{composite}.
\end{definitionbox}

\begin{workedbox}{Proof pattern: divisibility and linear combinations}
\textbf{Unpack:} assume $d\mid a$ and $d\mid b$, so $a=dr$ and $b=ds$ for
some $r,s\in\Z$. \textbf{Substitute and factor:} for any $m,n\in\Z$,
\[
ma+nb=m(dr)+n(ds)=d(mr+ns).
\]
\textbf{Close:} $mr+ns\in\Z$, so the display gives the required witness and
$d\mid(ma+nb)$.
\end{workedbox}

\begin{workedbox}{Proof pattern: a modulus that factors}
\textbf{Unpack:} suppose $p,q\ge2$, $d=pq$, and $x\equiv y\pmod d$. Then
$d\mid(x-y)$, so $x-y=kd=kpq$ for some $k\in\Z$. \textbf{Refactor:}
\[
x-y=p(kq)\quad\text{and}\quad x-y=q(kp).
\]
\textbf{Close:} $kq,kp\in\Z$, so $p\mid(x-y)$ and $q\mid(x-y)$, giving
\[
x\equiv y\pmod p,\qquad x\equiv y\pmod q.
\]
This direction needs only positive moduli with $p\mid d$ and $q\mid d$;
primality is not required. For instance, $103-4=99$ proves the same claim
from modulus $33$ to moduli $3$ and $11$.
\end{workedbox}

\columnbreak

\begin{methodbox}{A reliable proof skeleton}
\begin{enumerate}
  \item Translate each hypothesis using its definition.
  \item Name the integer witnesses.
  \item Manipulate the resulting equations.
  \item Exhibit the witness required by the conclusion.
  \item State why the witness is an integer.
\end{enumerate}
\end{methodbox}

\begin{workedbox}{Proof pattern: products preserve congruence}
\textbf{Unpack:} assume $a\equiv b\pmod m$ and $c\equiv d\pmod m$. Then
$a=b+mk$ and $c=d+m\ell$ for some $k,\ell\in\Z$. \textbf{Expand and factor:}
\[
ac-bd=m(b\ell+dk+mk\ell).
\]
\textbf{Close:} the parenthesized expression is an integer, so
$m\mid(ac-bd)$ and $ac\equiv bd\pmod m$.
\end{workedbox}

\begin{workedbox}{Proof pattern: odd squares modulo $8$}
\textbf{Represent the odd integer:} let $n=2k+1$. \textbf{Expand:}
\[
n^2-1=4k(k+1).
\]
\textbf{Find the extra factor of $2$:} one of $k,k+1$ is even, so
$k(k+1)=2t$ for some $t\in\Z$. Hence $n^2-1=8t$. Thus
$8\mid(n^2-1)$, proving $n^2\equiv1\pmod8$.
\end{workedbox}

\begin{workedbox}{Two short divisibility arguments}
\textbf{Transitivity:} if $a\mid b$ and $b\mid c$, write $b=ar$ and $c=bs$.
Substitution gives $c=(ar)s=a(rs)$; since $rs\in\Z$, $a\mid c$.

\textbf{Three consecutive integers:} let $P=n(n+1)(n+2)$. One factor is
even, so $P=2r$ for some $r\in\Z$. One factor is divisible by $3$, so
$P=3s$ for some $s\in\Z$. Combine these witnesses:
\[
P=3P-2P=3(2r)-2(3s)=6(r-s).
\]
Since $r-s\in\Z$, the definition gives $6\mid P$.
\end{workedbox}

\begin{warningbox}{Claims that need an extra hypothesis}
From $a\mid bc$, it does \emph{not} follow in general that $a\mid b$ or
$a\mid c$: $6\mid2\cdot3$, but $6$ divides neither factor. The implication
is valid when $a$ is prime (Euclid's lemma): if prime $p\mid bc$, then
$p\mid b$ or $p\mid c$.
\end{warningbox}

\begin{checklistbox}{Core readiness and format}
\textbf{A:} canonical division and congruence. \textbf{B:} base conversion.
\textbf{C:} modular arithmetic, cycles, squaring, and congruence proofs.
\textbf{D:} Euclidean algorithm and B\'ezout. \textbf{E:} direct proofs and
counterexamples. A live Core designates one problem from each section A--E
and lasts 20 minutes. Four fully correct problems earns \textbf{Meets}; all
five earns \textbf{Exceeds}. Every retake is a complete new form. Show enough
work for each answer and proof to be checked; approved accommodations apply.
\end{checklistbox}
\end{multicols}
\endgroup

\thispagestyle{fancy}
\end{document}
